\subsection{A compatible mixed-prime refinement}

\begin{statement}{D.28}{Proved}{An extra CM action through a mixed two--three kernel}
Keep $B=C_2*C_3$, $q:B\twoheadrightarrow Q$, $N$, $G_3$,
$H_2$ and $\mathcal P_3$ from D.25--D.27. Let
$V=\mathbb F_4$, viewed additively as $\mathbb F_2^2$.
The $Q$-action on $V$ factors through
\[
 Q\longrightarrow \mathrm{SL}_2(\mathbb F_3)/Q_8\simeq C_3;
 \qquad s|_V=I_2,\qquad
 r|_V=\begin{pmatrix}0&1\\1&1\end{pmatrix}.              \tag{D.28.1}
\]
In the ordered rank 97 graph basis of D.25, define
$\varphi:N^{\rm ab}/2\twoheadrightarrow V$ by the two
binary row masks
\[
 \texttt{1860e0d6cb6b6efc923133e5a}_{16},\qquad
 \texttt{1231c145bdbd87ed2d1f1a136}_{16},                 \tag{D.28.2}
\]
where bit $i$ is the coefficient of basis vector $i$, starting
at $i=0$. This is a $Q$-equivariant epimorphism. Its kernel is
independent of multiplying $\varphi$ by a nonzero scalar in
$\mathbb F_4$. Put
\[
 L=\ker(N\longrightarrow V),\quad H=B/L,\quad
 J=B/\bigl([L,L]L^3\bigr).                              \tag{D.28.3}
\]
These are specified discrete finite quotients, with
\[
 |H|=2304,\quad \operatorname{rank}L=385,\qquad
 1\longrightarrow\mathbb F_3^{385}\longrightarrow J
       \longrightarrow H\longrightarrow1.             \tag{D.28.4}
\]
The extra matrix pair
\[
 A=I,\qquad \beta=\begin{pmatrix}9&8\\4&9\end{pmatrix}\pmod{12}
                                                               \tag{D.28.5}
\]
admits a cusp-preserving automorphism of $J$ with multiplier
$\lambda=1$. It can be chosen to agree with D.27's alternate
action on the exact common quotient
\[
                    C=H\times_QG_3.                   \tag{D.28.6}
\]
Consequently there is one compatible automorphism, with
specified elliptic conjugators and multiplier one, of
\[
 \mathcal R=J\times_C
          \bigl(H_2\times_Q\mathcal P_3\bigr).          \tag{D.28.7}
\]
The canonical continuous map $\widehat B\twoheadrightarrow
\mathcal R$ is onto. One global cusp gauge gives compatible
nonarithmetical GT$_0$ shadows on every finite characteristic
refinement of its $G_6$ quotient within $\mathcal R$.

This refinement retains a mixed-prime action in the kernel:
\[
 \ker(J\to Q)\simeq
        \mathbb F_3^{385}\rtimes C_2^2,\qquad
 \dim_{\mathbb F_3}(\mathbb F_3^{385})^{C_2^2}=97.       \tag{D.28.8}
\]
The semidirect-product identification is noncanonical; the
kernel, its normal three-subgroup and their action are
canonical in (D.28.3). In particular $J$ is not a quotient
of any product over $Q$ of relative pro-prime completions
whose kernel is a direct product of pro-$p$ groups. The
maximal quotient of $J$ with nilpotent kernel over $Q$ is
exactly $C$. The full profinite realization of (D.28.5)
remains the all-finer-quotient question of D.22.
\end{statement}
\begin{proof}
Write $C_N=H_1(N\backslash\mathcal T,\mathbb F_2)$, where
$\mathcal T$ is the Bass--Serre tree used in D.25. A group
of order three acts irreducibly on the two-dimensional
space in (D.28.1), and
$\operatorname{End}_Q(V)=\mathbb F_4$. The exact equations
for a matrix $Z:C_N\to V$ are
\[
 ZM(s)=Z,\qquad ZM(r)=
       \begin{pmatrix}0&1\\1&1\end{pmatrix}Z.           \tag{D.28.9}
\]
The program \texttt{a\_mixed\_module\_explore.py} solves
these equations over $\mathbb F_2$. Their solution space
has dimension two, and the first basis matrix is (D.28.2).
Both displayed equations are checked on all 97 columns.
The matrix has rank two, so it is onto. Every nonzero
solution is obtained from it by a nonzero scalar in its
endomorphism field: this gives the stated independence
of its kernel. It also proves that $L$ is preserved by
any automorphism of $\widehat B$ inducing the identity on
$Q$, through the induced $Q$-equivariant action on
$N^{\rm ab}/2$. Normality of $L$ in $B$ follows directly
from (D.28.9).

For clarity, the finite extension $H$ is computed with its
extension class, not replaced by a split extension. Let
$P(q)$ be the fixed path chain to the representative of
$q\in Q$ in D.25. For $q,q'\in Q$ the chain
$P(q)+qP(q')-P(qq')$ is a cycle. Set
\[
 c(q,q')=\varphi\bigl(P(q)+qP(q')-P(qq')\bigr)\in V.
\]
Then $H$ consists of the 2304 pairs $(q,v)$ with law
\[
 (q,v)(q',v')=(qq',v+qv'+c(q,q')).                      \tag{D.28.10}
\]
The cocycle identity follows by cancelling the four path
expressions in the associativity equation. This is the
actual quotient $B/L$: the path construction sends its
generators to $(s,0),(r,0)$ and restricts to $\varphi$ on
$N$. Its kernel is therefore exactly $L$. The code also
checks that these generators reach all 2304 pairs.
The $H$-cover of the graph has 2304 edges and
$2304/2+2304/3=1920$ vertices. Its free fundamental group
has rank $2304-1920+1=385$, proving (D.28.4) after
abelianization modulo three.

Here are the complete finite lifting equations. Put
$C_L=H_1(L\backslash\mathcal T,\mathbb F_3)$ and use the
oriented graph's spanning-tree basis. Let $a,b\in H$ be
candidate reductions of $A,\beta$, and put
\[
 s'=bsb^{-1},\quad r'=ara^{-1},\quad t=sr.
\]
The finite checks first require $s'r'=t$ and that the
assignment $s\mapsto s'$, $r\mapsto r'$ is an automorphism
$\alpha_H$ of $H$. They evaluate its map on every one of
the 2304 elements and verify both generator transitions.
If $a_1,b_1\in C_L$ are the left correction coordinates
of the two conjugators, their cusp equation is
\[
 (1-s')b_1+s'(1-r')a_1=-d,                              \tag{D.28.11}
\]
where the actual path-chain defect is
\[
 d=(1-s')P(b)+s'\bigl((1-r')P(a)+aP(r)\bigr)-P(t).
                                                               \tag{D.28.12}
\]
All actions in these formulas are the left permutation
action on the edges of the $H$-cover. Thus both sides
are elements of the specified385-dimensional cycle
space; membership is checked by reconstruction from
their chord coordinates.

The compatibility conditions determine which finite
solution we use. In the coordinates (D.28.10), the
corrected conjugators of D.27 reduce to
\[
 a=(1,(1,0)),\qquad b=(q(\beta),(1,1)).                 \tag{D.28.13}
\]
To obtain the second coordinate one must include D.27's
right $s$-centralizer correction. The program evaluates
(D.28.2) on its full corrected homology mask, and
obtains the bit values1 and 3 in (D.28.13).
The desired reductions in $G_3$ are $A_3=1$ and
\[
                  P(\beta_3)=P(q(\beta))+2c_{39},     \tag{D.28.14}
\]
where $c_{39}$ is the basis cycle numbered39 in D.25.
These are exactly D.27's retained mod-three conjugators.

The inclusion $L\hookrightarrow N$ induces the map
$C_L\to N^{\rm ab}/3$ obtained by summing each edge chain
over the four edges above each edge of the $Q$-cover.
It has rank 97. Applying this explicit map to the two
corrected path chains gives194 further linear conditions
imposing (D.28.14). Together with (D.28.11) the resulting
system has rank 482 in 770 variables, and hence an affine
solution space of dimension 288. No arithmetic or
profinite lifting assumption enters this finite system.

The file \texttt{a\_mixed\_prime\_compatible.json}
contains its two complete385-coordinate solution vectors.
The 22nd tested vector in the deterministic enumeration
has invertible kernel action. The test of invertibility
uses the actual equivariant graph map: the base edge is
sent to the path
\[
        e_a+P(a)+a_1-P(b)-b_1,                          \tag{D.28.15}
\]
and edge $h$ is sent to its translate by $\alpha_H(h)$.
The induced385 by 385 matrix has rank 385 over
$\mathbb F_3$. The program additionally reconstructs all
385 free basis words of $L$, checks that each gives its
claimed fundamental cycle, and evaluates each basis
image from the substituted $s,r$ words. All 385 columns
agree with the matrix from (D.28.15). Finally the
corrected54658-letter cusp word
\[
       \beta s\beta^{-1}ArA^{-1}(sr)^{-1}
\]
is evaluated directly: its $H$-coordinate and all 2304
edge-chain coordinates are zero. These checks certify
existence independently of the random search method
used to locate the recorded vectors.

The resulting homomorphism $B\to J$ preserves $L$ because
its reduction is $\alpha_H$. It kills $[L,L]L^3$, and so
induces an endomorphism of $J$. Its action on the kernel
$L/[L,L]L^3$ is the invertible matrix just verified, and
its action on $H$ is an automorphism. It is therefore
surjective on $J$, hence an automorphism of the finite
group. The conjugators lie in the image of
$F=\ker(B\to S_3)$, since their images in $Q$ do; every
kernel correction lies in $L\subset F$. The multiplier
is exactly one.

We prove the common-quotient assertion, which is needed
to combine the witnesses. Write $M=L^{\rm ab}/3$ and
$V=N/L$. The kernel of $J\to Q$ is an extension
$1\to M\to\overline N\to V\to1$. Since $|V|=4$ is
invertible in $\mathbb F_3$, its extension class vanishes.
Explicitly, for a normalized additive two-cocycle
$b(v,w)$, put $h(v)=\frac14\sum_t b(v,t)$. Summing its
cocycle identity over $t\in V$ gives
$b(v,w)=vh(w)+h(v)-h(vw)$, so changing the section by
$-h$ splits the extension. This proves the semidirect
product in (D.28.8), with its stated ambiguity.

The finite covering map $L\backslash\mathcal T\to
N\backslash\mathcal T$ supplies a transfer. On homology
its composite with projection is multiplication by four;
averaging is an exact projection onto invariants because
four is invertible. More directly, the cellular chain
complex of the quotient graph is the coinvariant complex
of the covering graph. The coinvariant functor is exact
by averaging, so taking its first homology gives
\[
                 M_V\xrightarrow{\sim}N^{\rm ab}/3.
                                                               \tag{D.28.16}
\]
The explicit edge-sum map above is this isomorphism on
coinvariants. Consequently $\dim M^V=\dim M_V=97$,
whereas $\dim M=385$; the nontrivial summand has dimension
288. These dimensions and nontriviality are independently
checked using the two matrices of $V$ on $C_L$.

In a nilpotent finite quotient, the images of the two-
and three-Sylow subgroups commute. Thus every quotient
of $\overline N$ which is nilpotent kills $[V,M]$.
Conversely
$\overline N/[V,M]\simeq V\times M_V$ is nilpotent.
The maps $J\to H$ and $J\to G_3$ are defined because
$[L,L]L^3\subset L$ and $[L,L]L^3\subset[N,N]N^3$.
Their joint map is onto $H\times_QG_3$, by coprimality
of its two kernel orders. Equation (D.28.16) identifies
its kernel with $[V,M]$, proving (D.28.6) as the maximal
quotient of $J$ with nilpotent kernel over $Q$.

The kernel over $Q$ of
$H_2\times_Q\mathcal P_3$ is the direct product of a
finite two-group and a pro-three group. Every finite
quotient of that kernel is nilpotent. Its maximal common
quotient with $J$ is therefore at most $C$. Both factors
map onto $C$, so it is exactly $C$. Apply the elementary
subdirect-product description to the closed image of
$\widehat B$ in their product: its associated common
quotient is the quotient by the product of their kernels.
The equality just proved identifies that quotient with
$C$, and gives the surjection in (D.28.7).
The matrices and conjugators (D.28.13)--(D.28.14) show
that the two automorphisms agree on $C$. They consequently
define an automorphism of the indicated fibre product.

Choose one lift of its combined $\beta$ to $\widehat F$.
The cusp gauge of D.25, applied once to this lift, gives
a commutator representative and both hexagons throughout
the compatible inverse system. The CM pair is unchanged
under that gauge. Its two actions on the common CM field
in D.25 disagree, so none of these finite shadows retaining
that pair is induced by a Galois element. The quotient
(D.28.7) and its surjection from $\widehat B$ are the
precise source and target of the construction; no lift
of this quotient automorphism to $\widehat B$ is inferred.
\end{proof}
